\section{Proofs about Section~\ref{s:cbnNEW}}
%\section{Appendix LMCS for call-by-name}
\label{a:cbn}
             
\subsection{Some auxilary results}

% In the remainder we abbreviate $\encoNa M p$ as $\encoNa Mp$.
% \dsJ{i'd do so only in appendixes}


             

For $n>0$ we  define, for \piI\ abstractions  $A_1 \ldots A_n$:
\[\begin{array}{l}
\OUTU n {r_n}{r_0} {A_1 \ldots A_n}{} \defi 
\begin{array}[t]{l}
\res{r_1, \ldots r_{n-1}} \\
\! \! (\inpi r 0 {v_1} . \bopi v 1{x_{1},r_{1}}.    
 | \\
\inpi r 1 {v_2} . \bopi v 2{x_{2},r_{2}}    .  
 | \\
\ldots  \\
\inpi r {n-1} {v_n} . \bopi v n{x_{n},r_{n}}    .  
 \\
  !  \inp {x_1} {q} . {\app{A_1}{q} } | \\
  !  \inp {x_2} {q} . {\app{A_2}{q} } | \\
\ldots  \\
  !  \inp {x_n} {q} . {\app{A_n}{q} } )
  \end{array} 
  \end{array} 
\]



\begin{lemma} 
We have: 
\[
\begin{array}[t]{rcl}
\encoNa{ x M_1 \ldots M_n} {r_n}
\arrPI{\bout x{r_0}}
&\equiv  &
\begin{array}[t]{l}
\res{r_1, \ldots r_{n-1}} \\
\! \! (\inpi r 0 {v_1} . \bopi v 1{x_{1},r'_{1}}    .  \link {r'_{1}}{r_{1}}
 | \\
\inpi r 1 {v_2} . \bopi v 2{x_{2},r'_{2}}    .  \link {r'_{2}}{r_{2}}
 | \\
\ldots  \\
\inpi r {n-1} {v_n} . \bopi v n{x_{n},r'_{n}}    .  \link
{r'_{n}}{r_{n}} \\
  !  \inp {x_1} {q} . {\encoNa{M_1} {q} } | \\
  !  \inp {x_2} {q} . {\encoNa{M_2}{q} } | \\
\ldots  \\
  !  \inp {x_n} {q} . {\encoNa{M_n}{q} } )
\end{array}
\\ \ \\
&\wbPI &
\OUTU n {r_n}{r_0} {\encoN{M_1} \ldots \encoN{M_n}}{} 
\end{array}
\]
 \end{lemma}
\noindent
The two lemmas below are known results about encoding of call-by-name $\lambda$-calculus
into \piI\  \cite{San95i}. 


\begin{lemma} %[from  \cite{San95i}]
\label{l:sk}
 Let $M $ be a  $\lambda$-term.
$ $ 
\begin{enumerate}
\item  If  $a,b$ and $c$  are distinct names of the same type,  then 
  $\res b ( \link ab | \link bc ) \wb \link ac$; 


\item  If $x$ and $y$ are distinct  trigger  names and $y$ is not free in
$M$, then
  $\res{ x} (\link  xy  | \encoNa M r ) \wb \encoNa{ M
\sub y x} p$;   


\item  If $p$ and $r$ are distinct  continuation names, then 
 $\res{ r} (\link  r p | \encoNa M r ) \wb \encoNa  M p$.
%\iffull \else \sq \fi

\end{enumerate} 
\end{lemma} 


\begin{lemma} %[from  \cite{San95i}]
\label{l:beta_cbn}
If $M \longrightarrow M'$ then $\pairP {\encoNa M p } \longrightarrow \: {\contr}\; \pairP
{\encoNa        {M'} p }$
\end{lemma} 



\subsection{Operational correspondence}
\label{ss:oc_cbn}

We report here some results concerning the operational correspondence between 
OGS  and \piI terms. 


\begin{lemma}[Actions from encodings of 
$\LaoP$ terms]
\label{l:opt_op_cbn}
For any
  $\LL\in\LaoP$  and $\phi $ with 
 $p,\fv \LL \subseteq \phi$,  term  $\encoSPN L p$ has exactly
  one immediate transition, and exactly one of the following clauses holds:
\begin{enumerate}
\item $\encoSPN \LL p\arr{\bout pv} P $  
and $\LL$ is a value (hence of the form either  $ \lambda x. M $ or
$v$ ),  
 with $P =\encPPN { \pmap{v}{L}}{\phi \uplus \{v\}}$;

\item  
 $\encoSPN \LL p\arr{\bout x {q}}P$  
and $L =  
{E [x]}$, for some $E$, and 
$$P \contr  
  \encPPN { \pmap q{(E,p)}} {\phi \uplus \{q\}}
$$;

\item 
 $\encoSPN \LL p\arr{\bout v {x,r}}P$  
and $L =  
{E [v M]}$, for some $E,M$ 
and 
$$P \equiv 
  \encPPN { \pmap r{(E,p)}\cdot  \pmap x M } {\phi \uplus \{x,r\}}
$$;


\item 
$\encoSPN M p\arr\tau P$ and there is $N$  with $M\red N$ and
  $$P\exn \encoSPN  N p.$$ 
\end{enumerate}
\end{lemma}

\ifcutCBNLMCS 
\begin{proof}
\dsJ{I have checked (and thus rectified the def of $ \pmap q{(E,p)}$ in the encoding )} 
\end{proof}
\fi


\begin{lemma}[Actions from \piI encodings of  configurations]
\label{l:opt_con_cbn}
If  $\encoN{\conf {\gamma, \phi }} \arrPI \mu P$,
 then we have three possibilities:
\begin{enumerate}

\item $\mu $ is an input $v(x,q)$ and $\gamma 
= \pmap{v}{V}\cdot\gamma' $, and 
$$P \wb 
\encoN {
\conf{ V x,  q, \gamma', \phi \uplus \{ q,x\}}
} ;$$ 


\item $\mu $ is an input $x(q)$ and $\gamma 
= \pmap{x}{M}\cdot\gamma' $, and 
$$P  =
\encoN {
\conf{ M,  q, \gamma, \phi \uplus \{ q\}}
} ;$$ 


\item $\mu $ is an input $p(v)$ and $\gamma 
= \pmap{p}{E,q}\cdot\gamma' $, and 
$$P  = 
\encoN {
\conf{ E[v],  q, \gamma, \phi \uplus \{ v\}}
} .$$ 
\end{enumerate} 
\end{lemma} 


\begin{lemma}[from OGS to \piI, strong transitions]
\label{l:opcorrCON_cbn}
$ $ \begin{enumerate}

\item \caserule{IOQ}  If $F \arr{\initName (p)} F'$ then 
$\encoN{F } \arr  {\initName (p)}  
 \encoN{F'}$ ; 


\item  \caserule{Int} If $F \longrightarrow F'$ then 
$ \encoN F \longrightarrow  \contr \encoN {F'}$; 

\item  \caserule{PA} If $F \arr{\bout pv} F'$ then 
$ \encoN F \arr{\bout pv}  \equiv  \encoN {F'}$; 
\item  \caserule{PVQ} If $F \arr{\bout v {y,q}} F'$ then 
$ \encoN F \arr{\bout v {y,q}}   \equiv  \encoN {F'}$; 
\item  \caserule{PTQ} If $F \arr{\bout x {q}} F'$ then 
$ \encoN F \arr{\bout x {q}}  \contr  \encoN {F'}$; 
\item  \caserule{OA} If $F \arr{\inp  {q} v} F'$ then 
$ \encoN F\arr{\inp  {q} v}    \encoN {F'}$; 
\item  \caserule{OVQ} If $F \arr{\inp   v{y,q}} F'$ then 
$ \encoN F\arr{\inp   v{y,q}}  \wb  \encoN {F'}$; 
\item  \caserule{OTQ} If $F \arr{\inp   x{q}} F'$ then 
$ \encoN F\arr{\inp   x{q}}    \encoN {F'}$.
\end{enumerate}
\end{lemma} 

% We can try to  abbreviate the above as follows: 

% \begin{theorem}[from OGS to \piI, strong transitions]
% \label{t:opcorrCON_cbnshort}
% $ $ \begin{enumerate}

% \item   If $F \arr{\initName (p)} F'$ then 
% $\encoN{F } \arr  {\initName (p)}  
%  \encoN{F'}$ ; 

% \item  If $F \longrightarrow F'$ then 
% $ \encoN F \longrightarrow  \contr \encoN {F'}$; 

% \item  If $F \arr{\mu} F'$ where $\mu$ is an input or an output action
% % not $\tau$ or of the form 
% % $\initName (p)$ 
% then 
% $ \encoN F \arr{\mu}  \wb  \encoN  {F'}$.
% \end{enumerate}
% \end{theorem} 

% \finish{if useful we could use $\ell$ to range over the actions required in (3) above} 


\begin{theorem}[from OGS to \piI, traces]
\label{t:tracesGSpi_cbn}
If $F$  has the trace $s$ then also $ \encoN F$ has the trace $s$.
\end{theorem} 

\begin{proof}
Straightforward induction on the length of the trace. 
\end{proof} 
\noindent
Now the converse direction. 


\begin{lemma}[from \piI to OGS, strong transitions]
\label{l:opcorr_cbnshort_piI}
$ $ \begin{enumerate}

\item   
If $\encoN{F } \arr  {\initName (p)}  
P$, 
then there is $F'$ such that 
 $F \arr{\initName (p)} F'$
 and
 $P = 
 \encoN{F'}$ ; 

\item  If 
$ \encoN F \longrightarrow P$, then 
then there is $F'$ such that 
$F \longrightarrow F'$ and  
$ P  \contr \encoN {F'}$; 

\item  If
$ \encoN F \arr{\mu} P$, 
where $\mu$ is an output,
 then 
then there is $F'$ such that 
$  F \arr{\mu} F' $ and 
$P  \contr  \encoN  {F'}$; 

\item  If
$ \encoN F \arr{\mu} P$, 
where $\mu$ is an input,  and 
$ \encoN F$ cannot  perform an output or a $\tau$,  
then there is $F'$ such that 
$  F \arr{\mu} F' $ and 
$P  \wb  \encoN  {F'}$. 
\end{enumerate}
\end{lemma} 
\noindent
In the modified transition system $\arrN {\mu'}$  for \piI, where output and silent transitions have
priority over input transitions, we can then derive from the above results the  full
abtractions on traces, reported in Corollary~\ref{c:CBNtraces}. 

\subsection{Concurrent Operational Game Semantics}


We have reported in the main text the operational correspondence result for strong
transitions. We can thus derive the 
operational correspondence result for weak
transitions.
 For silent transition, we proceed by
induction on the number of strong transitions involved. 

\begin{theorem}[from \COGS to \piI, weak transitions]
\label{t:CBNpogs_short_weak}
$ $ \begin{enumerate}

\item  If $F \Longrightarrowp F'$ then 
$ \encoSN F \LongrightarrowPI \; \contr \encoSN {F'}$; 

\item  If $F \ArrC{\ell} F'$  then 
$ \encoSN F \ArrPI{\ell} \;  \wbPI  \encoSN  {F'}$.
\end{enumerate}
\end{theorem} 


\begin{theorem}[from  \piI to \COGS, weak transitions]
\label{t:CBNpogs_short_converse_weak}
$ $ \begin{enumerate}

\item  If
$ \encoSN F \LongrightarrowPI  P$, then there is $F'$ s.t.\  
 $F \Longrightarrowp F'$ and 
$P 
\contr \encoSN {F'}$;

\item  If
$ \encoSN F \ArrPI{\ell}  P$, then there is $F'$ s.t.\  
 $F \ArrC{\ell} F'$ and 
$P 
\wbPI \encoSN {F'}$.
\end{enumerate}
\end{theorem} 



\begin{corollary}
\label{c:CBNpogs_pi}
 $F$ and $ \encoSN F$ have the same set of traces. 
\end{corollary}  

\begin{proof}
By induction on the length of a trace performed by 
 $F$  or $ \encoSN F$: thus one shows that, for a trace $t$,
\begin{itemize}
\item  If $F \ArrC{t} F'$  then 
$ \encoSN F \ArrPI{t}\;  \wbPI  \encoSN  {F'}$;
\item  If
$ \encoSN F \ArrPI{t}  P$, then there is $F'$ s.t.\  
 $F \ArrC{t} F'$ and 
$P 
\wbPI \encoSN {F'}$.\qedhere
\end{itemize} 
\end{proof} 



\begin{corollary}
\label{c:CBNpogs_pi_traces}
We have: 
 $\FF_1 \TEp \FF_2$ iff 
 $ \encoSN {\FF_1} \TEPI  \encoSN {\FF_2}$.  
\end{corollary}  


\begin{corollary}
\label{c:CBNpogs_pi_bis}
we have: 
 $\FF_1 \wbC \FF_2$ iff 
 $ \encoSN {\FF_1} \wbPI  \encoSN {\FF_2}$.  
\end{corollary}  

\ifcutCBNLMCS
\begin{proof}
In one direction, one shows that 
\[  \{ (    \encoSN {\FF_1},  \encoSN {\FF_2} ) \st \FF_1 \wbC \FF_2 \}\]
is a weak bisimulation up-to $\wb$. 
For the other direction, one shows that 
\[  \{ (    {\FF_1},  {\FF_2} ) \st \encoSN {\FF_1} \wbPI   \encoSN {\FF_2}\}\]
is a bisimulation.
\end{proof} 
\fi

\noindent
A configuration $F$ \emph{is a singleton} if its domain has only one
element.



\begin{lemma}
\label{l:CBNtr_bisi_sing}
For any singleton $\FF,\GG$, we have: $
{\encoSN \FF} \TEPI {\encoSN \GG}
$ iff 
${\encoSN \FF} \wb {\encoSN \GG}$.
\end{lemma} 

\begin{proof}
Weak bisimilarity implies trace equivalence, hence the direction from right to left. 

For the converse, we show
 that 
the relation 
\[ 
(\encoSN F,  \encoSN G) \st \mbox{ $F,H$ are singleton with $F \TEp
G$}
\]
is a bisimulation up-to contexts and expansion. 
The proof uses Lemmas~\ref{l:immediate_tr} and \ref{l:deco}. 
\finish{why this holds: the operational correspondence resuts; the fact that the names in
  the domain of a configuration do not appear anywhere else in the configuration} 
\end{proof} 
\noindent
We think that the above lemma can be generalised to aribitrary configurations; the current
statement is however sufficient for our purposes.
By combining Lemma~\ref{l:CBNtr_bisi_sing} and 
Corollaries~\ref{c:CBNpogs_pi_traces} and 
\ref{c:CBNpogs_pi_bis} we then  derive  Theorem~\ref{t:CBNfa_all}. 

\noindent
The following result holds for call-by-name but not for call-by-value.
\finish{above: check again...  is it so, really? } 
\finish{the result can be generalised to arbitrary configurations. It needs a bit more of
  work. Do we want to do it? } 
\begin{lemma}
\label{l:CBNonly}
Suppose $\gamma$ and $\delta$ (resp.\   
$\gamma'$ and $\delta'$) have the same domains.
Then 
 $\conf{\gamma \cdot \gamma',\phi } \TEp 
\conf{ \delta  \cdot \delta',\phi } $
implies
 $\conf{\gamma,\phi } \TEp 
\conf{ \delta  ,\phi } $
 \end{lemma}

\begin{proof}
This is a consequence of the fact that any action from the components $\gamma', \delta'$ 
has, as subject, a name in the domain of these components. Such a name does not appear in $
\gamma , \delta $, and thus it will not appear in any action of their traces.  
\end{proof} 

\ifcutCBNLMCS
\subsection{Up-to techniques for games}
\label{ss:CBNupto_games}
\finish{again, this section is essentially the same as fo call-by-value, with some
  notations revisited} 

The goal of this section is to trasnsfer up-to techniques from 
\piI  onto GS. We will then use such techniques to prove that
\COGS and ordinary OGS yield the same semantics on
$\lambda$-terms. 

 

A relation $\R$ on configuration is \emph{well-formed} 
if  the domains of any two configurations 
$\FF_1$ and $ \FF_2$
with  $\FF_1 \RR \FF_2$  are the same.

Two configurations are \emph{disjoint} if the set of names that appear
in their flat  domain are disjoint. 

If 
$ \FF = \conf{ A , \gamma,\phi }$
and 
$ \GG = \conf{ B , \delta ,\phi' }$
are disjoint, then 
\[ \appendCon \FF \GG  \defi
 \conf{ A \cdot B , \gamma \cdot \delta,\phi\uplus \phi' }
 \] 
\finish{above: should $\phi'$ be $\phi$?} 

Below, all relations on configurations are meant to be well formed. 

Given a well-formed relation $\R$ we write
$\appendRel \R $ for the relation
\[ 
\{ (\FF_1,\FF_2) \st
\mbox{there is $\GG$ s.t.\ }
 \FF_i = \appendCon{\FF'_i}\GG   \mbox{ $i=1,2$}  \mbox{ and }  \FF'_1\RR \FF'_2\}
\]


Finally, $\uptoComp \R $ is the reflexive and transitive closure of 
$\appendRel \R $

Note: equality on environments ignores the order of their components
(ie, it is  seen as a set).


\begin{definition}
\label{d:uptoComp-cbn}
A relation on configurations is a \emph{bisimulation up-to
  composition} if whenever $\FF_1 \RR \FF_2$: 
\begin{itemize}
\item if $\FF_1 \arrC\mu \FF_1'  $ then there is $\FF_2' $ such that 
$\FF_2 \arrC\mu \FF_2'  $ and 
$\FF'_1 \uptoComp \R \FF'_2$

\item  the converse
\end{itemize}  
\end{definition} 

If $R$ is a relation on configuration, then 
$\encoSN \R$ is the relation on \piI terms obtained by mapping each pair
in $\R$ in the expected manner: 
\[ \encoSN \R \defi \{ (\encoSN {\FF_1}, \encoSN {\FF_2}) \st \FF_1 \RR \FF_2 
 \} \] 


\begin{theorem}
\label{t:uptoComp-cbn}
If 
$\R$ is bisimulation up-to
  composition then $\R \subseteq {\wbC}$.
\end{theorem} 

\begin{proof}
By showing that 
$\encoSN \R$ is a bisimulation up-to expansion and contexts in 
\piI.
\finish{all is needed here is a result of operational correspondence} 
\finish{upto transitive closure is captured by having ``upto polyadic contexts''} 
\end{proof} 


\subsection{Relationship between \COGS and OGS}

\finish{actually here $M$ ranges over  extended $\lambda$-terms} 

For any term $M$, its traces in ordinary OGS are a subset of the traces in 
Concurrent OGS (\COGS), precisely, the traces with an alternation between player and
opponent moves. 
\finish{correct?}   
In other words, the  syntax of a trace of $M$   reveals whether that trace is also a trace
for $M$ in ordinary OGS.  
As a consequence, trace equivalence in \COGS implies trace equivalence in OGS. 

\begin{lemma}
\label{l:CBNPOGS_OGS}
If $\FF_1 \TEp \FF_2  $ then also $\FF_1 \TEPI \FF_2$. 
\end{lemma}


For the opposite direction, we exploit the 
`bisimulation up-to
  composition' technique. 


A configuration $\FF$ is a \emph{singleton} if the domain of $\FF$ has only 1 element. 
That is, $\FF$ is either of the form 
$\conf{\pmap{p}{M}, \emptyset, \phi }$, or 
$\conf{\pmap{v}{V}, \emptyset, \phi }$, or 
$\conf{\pmap{p}{(E,q)}, \emptyset, \phi}$, or
$\conf{\pmap{x}{M}, \emptyset, \phi}$.

\finish{whatch out for the initial term....}  


In the proof, 
we write $\FF \subseteq \GG$ when, as partial maps, $\GG$ is an extension of $\FF$. 
\finish{and moreover there is inclusion on the $\phi$ components} 
As for call-by-value, we distinguish between transitions that a configuration can make in \COGS and \AOGS, 
writing $\arrC \mu  $ for the former, and $\arrA\mu$ for the latter. 
Similarly for weak transitions $\ArrC \mu  $ and $\ArrA \mu  $.


\begin{lemma}
\label{l:CBNOGS_COGS}
If $\FF_1, \FF_2  $ are singleton and 
$\FF_1 \wbA \FF_2  $, then also $\FF_1 \wbC \FF_2$. 
\end{lemma}  

\begin{proof}
Let $\R $ be the relation on configurations with $\FF \RR \GG$ if:
\begin{enumerate}
\item $\FF,\GG$ are singleton, and
\item  either 
\begin{enumerate}
\item
there are passive  $\FF', \GG'$ with $\FF \subseteq \FF'$ and $\GG \subseteq \GG'$ and $\FF' \wbA
  \GG'$,
\item   or 
there are $p,M,N,  \gamma , \delta $ with
$\FF = \conf{\pmap{p}{M},\phi}$,  
$\GG = \conf{\pmap{p}{N},\phi}$, and 
$\conf{\pmap{p}{M}, \gamma,\phi } \wbA 
\conf{\pmap{p}{N}, \delta,\phi  } $
\end{enumerate} 
\end{enumerate}
\finish{initial terms are hopefully special cases of 2b...} 
\finish{i think one has also to impose that related singleton have the same flat domains} 


We show that $\R$ is  a bisimulation up-to
  composition.
We distinguish the cases in which (2.a) or (2.b) holds.


First suppose $\FF \RR \GG$ because (2.a) holds, and there are $\FF',\GG'$ as in (2.a).
Suppose 
\begin{equation}
\label{e:CBNKoaH}
 \FF \arrC  \mu \FF_1 
\end{equation} 
in \COGS.  The rule can be OA or OVQ, or OTQ.  We assume it is OA; the cases  for  OVQ and
OTQ are similar.  
We thus have, for some $p,q,E,v$ 
$$
\begin{array}{rcl}
 \FF &=&  \conf{ \gamma_1 ,\phi} \mbox{ with }  \gamma_1 =    \pmap{p}{(E,q)} \\
\mu & = & p(v) \\
\FF_1 & = & \conf{A, \gamma_1, \phi \uplus \{v\}} \mbox{ with }  A = \pmap{q}{E[v]}
\end{array} 
$$
Moreover,  we have $\GG = \conf{ \delta_1,\phi }$  with $  \delta_1 =    \pmap{p}{(\DD,q)}$, for
some $\DD$ (as related singleton configurations have the same flat domain), and, for
$\FF',\GG'$ as in (2.a), for some $\gamma, \delta $:  
$$
\begin{array}{rcl}
\FF' & = & \conf{ \gamma_1 \cdot \gamma  ,\phi} \\ 
\GG' & = & \conf{ \delta_1 \cdot \delta  ,\phi} 
\end{array} $$


We also have 
$ \FF' \arrA  \mu \conf{A,  \gamma_1 \cdot \gamma  ,\phi\uplus \{v\}} \defi \FF''$. 
Since $\FF' \wbA \GG'$, 
\[ \GG'  \arrA  \mu \conf{B,  \delta_1 \cdot \delta  ,\phi \uplus \{v\}}  \defi \GG''\] 
for 
$B \defi \pmap{q}{\DD[v]}$, 
with $\FF'' \wbA \GG''$.

Let $\GG_1 =   \conf{B, \delta_1}$.
We also have 
\[ \GG \arrC  \mu \GG_1 \] 
This is sufficient to match \reff{e:CBNKoaH}, as 
$\FF_1 \uptoComp \R  \GG_1$: indeed we have  
both $\conf{ \gamma_1,\phi\uplus \{v\}} \R  \conf{ \delta_1,\phi\uplus \{v\}}$ and  
(using clause (2.b))
$\conf{ A, \emptyset , \phi\uplus \{v\}} \R  \conf{ B, \emptyset, \phi\uplus \{v\}}$.

\finish{above : i am implicitly using several times a result of the kind 
$\conf{..., \phi} \wbA \conf{..., \phi}$ implies 
$\conf{..., \phi'} \wbA  \conf{..., \phi'}$ where $\phi'$ is larger than $\phi$ 
} 


Now the case (2.b);  let $p,  \gamma , \delta $ as in (2.b).
 There are 4 possibilities of transitions, corresponding to rules 
Int, PA, PVQ, and PTQ. The case of 
Int  is straightforward (also note that bisimilarity is preserved by internal
moves).  We only look at PVQ, as the others are simpler. 
Thus we have 
$\FF = \conf{\pmap{p}{E[v M_1]},\emptyset ,\phi }$, for some $E,v,M_1$,
and
$\GG = \conf{\pmap{p}{N]},\emptyset ,\phi }$
 (here again we use the property that  the flat domains of $\FF$ and
$\GG$ are the same) , 
 and then 
\begin{equation}
\label{e:CBNKpqH}
\FF  \arrC{\questPV{v}{y}{q}}
     \conf{ \pmap{y}{M_1}\cdot \pmap{q}{(E,p)},\emptyset , \phi\uplus\{y,q\}}  \defi \FF_1
 \end{equation} 
 We also have  
\[
\FF' \defi \conf{\pmap{p}{E[v M_1]}, \gamma,\phi }  
 \arrA{\questPV{v}{y}{q}}
     \conf{ \pmap{y}{M_1}\cdot \pmap{q}{(E,p)},\gamma, \phi\uplus\{y,q\}} 
\defi \FF''\]
Since 
$\FF' 
\wbA 
\conf{\pmap{p}{N}, \delta,\phi  } \defi \GG' $ 
we must have $N \LongrightarrowA  \DD[v N_1]$, for some $\DD,N_1$ so that 
we  can derive
\[
\GG' \defi \conf{\pmap{p}{\DD[v N_1]}, \delta,\phi }  
 \ArrA{\questPV{v}{y}{q}}
     \conf{ \pmap{y}{N_1}\cdot \pmap{q}{(\DD,p)},\delta, \phi\uplus\{y,q\}} 
\defi \GG''\]
with 
$\FF'' \wbA \GG''$. 

Let $$
\GG_1 \defi      \conf{ \pmap{y}{N_1}\cdot \pmap{q}{(\DD,p)},\emptyset , \phi\uplus\{y,q\}} 
.$$
We also have $\GG \ArrC
{\questPV{v}{y}{q}}
 \GG_1$. 

From $\FF'' \wbA \GG''$, appealing to (2.a)  we deduce that both 
$     \conf{ \pmap{y}{M_1},\emptyset , \phi\uplus\{y,q\}} \RR
     \conf{ \pmap{y}{N_1},\emptyset , \phi\uplus\{y,q\}}$ 
and 
$
     \conf{ \pmap{q}{(\DD,p),\emptyset , \phi\uplus\{y,q\}} 
 \RR
\pmap{q}{(E,p)},\emptyset , \phi\uplus\{y,q\}}$ hold.
Hence $\FF_1 \uptoComp \R  \GG_1$.

In summary, as an answer to the challenge 
\reff{e:CBNKpqH}, we have found  $\GG_1$ such that 
 $\GG \ArrC
{\questPV{x}{y}{q}}
 \GG_1$ and  
$\FF  _1 \uptoComp \R  \GG_1$; this closes the proof.
  \end{proof} 

\dsJ{note proof : done  assuming permanent environments, check again with ephemeral ones} 


We can thus derive 
Corollary~\ref{c:CBNOGS_COGS}, reported in the main text, and  stating that all forms of
trace equivalence and bisimilarity, on GS or \piI\ semantics, coincide.   
%  and  
% \begin{corollary}
% \label{c:CBNOGS_COGS}
% On $\lambda$-terms, relations $\wbA$, $\TEA$, 
% $\wbC$, $\TEp$  coincide. 
% \end{corollary} 
%
We remark that such a  result  cannot be extended to arbitrary configurations.
For instance, 
we have 
\[ \conf {\pmap  p \Omega ,  \pmap  x M , \phi} 
\TEA
\conf {\pmap  p \Omega ,  \pmap  x N , \phi} 
\]
for any $M,N$, which need not be the case for 
$\TEp$, as the visible behaviours of $M$ and $N$ may differ. 


However the result does  hold for passive configurations, or for configurations reachable from
$\lambda$-terms.


\finish{intuitively, why it works: the operational correpondence with pi; the fact that all new entries in the environment are on new names and the first action will be on such names. These properties hold both for CBN and for CBV} 

\fi
